% 363_Hodge_T4Z3_De_ch.tex
% Johann Pascher, 13. Sept. 2026
% Die Hodge-Vermutung und die T⁴/Z₃-Geometrie der FFGFT

\providecommand{\BM}{\textbf{[B]}}
\providecommand{\KM}{\textbf{[K]}}
\providecommand{\SM}{\textbf{[S]}}
\providecommand{\SetzM}{\textbf{[SETZUNG]}}
\providecommand{\QM}{\textbf{[Q]}}
\providecommand{\EM}{\textbf{[E]}}
\providecommand{\XM}{\textbf{[X]}}

\chapter*{Hodge-Theorie auf T\textsuperscript{4}/Z\textsubscript{3} im Rahmen der FFGFT}
\addcontentsline{toc}{chapter}{Hodge-Theorie auf T⁴/Z₃ im Rahmen der FFGFT}

\begin{abstract}
Das zweite Millennium-Problem der Mathematik — die Hodge-Vermutung — fragt,
ob jede Hodge-Klasse einer glatten projektiven Varietät eine rationale
Linearkombination algebraischer Zykel ist. Für die Basisgeometrie der FFGFT,
den Orbifold $T^4/\mathbb{Z}_3$, lässt sich diese Frage positiv beantworten:
Der komplexe 2-Torus $T^4$ ist eine abelsche Fläche, für die die
Hodge-Vermutung bekannt gilt; die neun $\mathbb{Z}_3$-Fixpunkte liefern
algebraische 0-Zykel; der Frobenius-Automorphismus $x \mapsto x^3$ auf
$GF(27)^*$ stimmt mit dem algebraischen Frobenius auf der étalen Kohomologie
überein; und die Sektorpaarung $k \leftrightarrow -k$ (Dok.~343) ist die
Hodge-Symmetrie $H^{p,q} \leftrightarrow H^{q,p}$. Das Millennium-Problem in
seiner allgemeinen Form — für beliebige glatte projektive Varietäten — bleibt
offen. Was FFGFT liefert, ist eine strukturelle Vollständigkeit für die
eigene Geometrie.
\end{abstract}

% ============================================================
\section*{Statusmarkierungen}
% ============================================================
\begin{center}
\begin{tabular}{@{}lp{10cm}@{}}
\textbf{[B]} & Algebraisch bewiesen. \\[3pt]
\textbf{[K]} & Numerisch bestätigt. \\[3pt]
\textbf{[S]} & Strukturell plausibel, physikalische Identifikation offen. \\[3pt]
\textbf{[E]} & Etablierte externe Mathematik. \\[3pt]
\textbf{[Q]} & Aus der Literatur übernommen, korrekt zitiert. \\[3pt]
\textbf{[SETZUNG]} & Axiomatische Setzung. \\
\end{tabular}
\end{center}

% ============================================================
\section{Das Millennium-Problem}
% ============================================================

Das Clay Mathematics Institute formuliert die Hodge-Vermutung wie folgt
\QM: Für eine glatte projektive algebraische Varietät $X$ über $\mathbb{C}$
ist jede Hodge-Klasse — also jede Kohomologieklasse vom Typ $(k,k)$ — eine
rationale Linearkombination von Klassen algebraischer Zykel der Kodimension $k$.

Präziser: Die Hodge-Zerlegung
\[
H^n(X,\mathbb{C}) = \bigoplus_{p+q=n} H^{p,q}(X)
\]
teilt die Kohomologie in Typen; die Hodge-Klassen sind die Elemente von
$H^{k,k}(X) \cap H^{2k}(X,\mathbb{Q})$. Die Vermutung besagt, dass diese
Klassen von echten geometrischen Objekten — Untervarietäten, deren
Kohomologieklassen rational kombiniert werden — erzeugt werden. Für glatte
projektive Kurven und Flächen vom Typ $(1,1)$ ist dies durch den
Lefschetz-Satz über $(1,1)$-Klassen bekannt \EM. In allgemeiner Dimension
ist die Vermutung offen.

% ============================================================
\section{Der komplexe 2-Torus als abelsche Fläche}
% ============================================================

Die reale Geometrie $T^4 = S^1 \times S^1 \times S^1 \times S^1$ trägt in
der FFGFT die Struktur eines \emph{komplexen 2-Torus}: zwei komplexe
Koordinatenpaare $(z_1, z_2) \in \mathbb{C}^2 / \Lambda$, wobei $\Lambda$
das Gitter der Modenzahlen ist \SetzM\ (Dok.~001).

Ein komplexer 2-Torus ist eine abelsche Fläche — eine abelsche Varietät der
komplexen Dimension~2. Für abelsche Varietäten ist die Hodge-Vermutung in
den Fällen bekannt, die hier relevant sind \EM\ \QM:
\begin{itemize}
\item Für abelsche Flächen (komplexe Dimension~2) gilt die Vermutung vollständig
      für alle Hodge-Klassen, da $H^{1,1}$ durch den Lefschetz-Satz abgedeckt
      ist und $H^{2,2}$ durch die Fundamentalklasse erzeugt wird.
\item Produkte von elliptischen Kurven und Tori über $\mathbb{Q}$ sind
      vollständig durch algebraische Zykel überspannt.
\end{itemize}

\noindent \BM\ \textit{Satz A:} $T^4$ als komplexer 2-Torus erfüllt die
Hodge-Vermutung: Alle $(k,k)$-Klassen sind rationale Linearkombinationen
algebraischer Zykel.

Beweis: Für $k=0$ und $k=4$: Trivial durch Fundamental- und Punktklasse.
Für $k=1$: Lefschetz-Satz auf $(1,1)$-Klassen (bekannter Satz \EM).
Für $k=2$: $H^{2,2}(T^4)$ wird erzeugt durch Produkte von $(1,1)$-Klassen,
also algebraisch. \hfill$\square$

% ============================================================
\section{Die neun Z₃-Fixpunkte als algebraische 0-Zykel}
% ============================================================

Unter der $\mathbb{Z}_3$-Wirkung auf $T^4$ gibt es genau neun Fixpunkte
(Dok.~001/330) \BM. Diese Fixpunkte sind algebraische 0-Zykel — Punkte auf
der Varietät $T^4/\mathbb{Z}_3$ mit wohldefinierten Kohomologieklassen.

\BM\ \textit{Satz B:} Die neun $\mathbb{Z}_3$-Fixpunkte von $T^4/\mathbb{Z}_3$
sind algebraische 0-Zykel der Kodimension~4 (als Punkte). Ihre Kohomologie-
klassen erzeugen den Twisted-Sektor der Orbifold-Kohomologie.

In der Orbifold-Kohomologie (Chen--Ruan) \QM\ zerfällt:
\[
H^*_{\mathrm{CR}}(T^4/\mathbb{Z}_3) = H^*(T^4)^{\mathbb{Z}_3}
\;\oplus\; \bigoplus_{\text{9 Fixpunkte}} H^*(\{p\})
\]
Der zweite Summand ist algebraisch: Jeder Punkt ist ein algebraischer
0-Zykel, seine Kohomologieklasse ist trivial algebraisch. \EM

Unter der Witt-Paar-Struktur von Dok.~341 sind:
\begin{itemize}
\item \textbf{Gerade Indizes 0,2,4,6} $=$ $\mathbb{Z}_3$-Fixpunkte
      $=$ algebraische 0-Zykel \BM
\item \textbf{Ungerade Indizes 1,3,5,7,8} $=$ Dreier-Orbits
      $=$ höherdimensionale algebraische Zykel \BM
\end{itemize}

Das ist buchstäblich die Hodge-Aussage: Topologische Klasse $=$ algebraischer
Zykel, für alle Klassen im Twisted-Sektor.

% ============================================================
\section{Frobenius-Automorphismus und algebraischer Frobenius}
% ============================================================

Dies ist die tiefste Verbindung zwischen FFGFT-Algebra und Hodge-Theorie.

\subsection{Der Galois-Frobenius auf GF(27)*}

In Dok.~339 wird bewiesen \BM: Der Frobenius-Automorphismus auf
$GF(27) = GF(3^3)$ ist die Abbildung $x \mapsto x^3$. Er wirkt auf
$GF(27)^* \cong \mathbb{Z}_{26}$ und erzeugt die Sektorstruktur:
Fixpunkte $\{+1,-1\}$ (massiver Sektor), acht Dreier-Orbits
(Gluonen/masseloser Sektor).

\subsection{Der algebraische Frobenius auf der étalen Kohomologie}

In der Zahlentheorie der algebraischen Geometrie \QM\ ist für eine Varietät
$X$ über $GF(q)$ der geometrische Frobenius die Abbildung
$F_q: x \mapsto x^q$ auf Punkten. Er induziert eine Abbildung
$F_q^*$ auf der $\ell$-adischen Kohomologie $H^i_{\mathrm{\acute{e}t}}(X,\mathbb{Q}_\ell)$.
Die Eigenwerte von $F_q^*$ sind algebraische Zahlen, die nach dem
Weil-Deligne-Theorem \EM\ die Hodge-Zahlen von $X$ bestimmen.

\BM\ \textit{Satz C:} Für $T^4/\mathbb{Z}_3$ über $GF(3)$ ist der
geometrische Frobenius $x \mapsto x^3$ auf $GF(27)^*$ identisch mit dem
algebraischen Frobenius auf $H^1_{\mathrm{\acute{e}t}}$.

Beweis (Skizze): $T^4$ über $GF(3)$ hat $GF(3)$-rationale Punkte, die durch
das Gitter $\Lambda$ parametrisiert sind. Der Frobenius $F_3$ wirkt als
$x \mapsto x^3$ auf den Koordinaten. Die Eigenwerte in $H^1$ sind die
26. Einheitswurzeln in $\overline{\mathbb{Q}}_\ell$ — exakt die Elemente von
$GF(27)^*$. Die Galois-Gruppe $\mathrm{Gal}(GF(27)/GF(3)) = \langle x \mapsto x^3 \rangle$
ist die Frob-Wirkung. \hfill$\square$

\SM\ \textit{Bemerkung:} Die physikalische Interpretation — Frobenius-Orbits
als Teilchensektoren — ist eine FFGFT-Einordnung und kein Teil des
mathematischen Beweises.

% ============================================================
\section{Sektorpaarung als Hodge-Symmetrie}
% ============================================================

In Dok.~343, Satz~D \BM\ wird bewiesen: Die Sektorpaarung
$k \leftrightarrow -k \pmod{13}$ auf den primitiven GF(27)*-Klassen paart
$f_1 \leftrightarrow f_2$ und $f_3 \leftrightarrow f_4$.

\BM\ \textit{Satz D:} Diese Sektorpaarung ist die algebraische Realisierung
der Hodge-Symmetrie $H^{p,q} \leftrightarrow H^{q,p}$ durch komplexe
Konjugation.

Beweis: Die Hodge-Zerlegung auf $T^4$ hat $H^{1,0}$ und $H^{0,1}$ als
konjugiertes Paar; ebenso $H^{2,1}$ und $H^{1,2}$. Die Konjugation
$\overline{\cdot}: H^{p,q} \to H^{q,p}$ ist das algebraische Datum, das
reale Kohomologieklassen auszeichnet. Auf $GF(27)^*$ ist die entsprechende
Abbildung die Inversion $g^k \mapsto g^{-k}$, also die Negation des
Exponenten in $\mathbb{Z}_{26}$. Sie ist zu unterscheiden von der
Körpernegation $e \mapsto -e = e\cdot g^{13}$: Diese schickt jeden
Ordnung-26-Orbit auf einen Ordnung-13-Orbit und ist daher keine Paarung
innerhalb der primitiven Klassen (Prüfskript pruef\_363b, Teil L3).
Die Inversion dagegen ist eine Involution auf den vier primitiven
Klassen ohne Fixorbit; sie reduziert $\{f_1,f_2,f_3,f_4\}$ auf zwei
reale Klassen — exakt die Hodge-Symmetrie. \hfill$\square$

Aus Dok.~343 Satz~D‴ folgt \BM: $\{f_3,f_4\}$ ist die Dirac-Schicht
(geladene Leptonen), $\{f_1,f_2\}$ die Majorana-Schicht (Neutrinos). Die
Hodge-Paarung trennt damit genau die physikalischen Sektoren.

% ============================================================
\section{Die χ-Klassen als algebraische Basis}
% ============================================================

In Dok.~342 \BM\ werden die acht irreduziblen kubischen Polynome über
$GF(3)$ klassifiziert:
\begin{itemize}
\item Vier primitive Polynome der Ordnung~26: $f_1, f_2, f_3, f_4$
\item Vier Polynome der Ordnung~13: $g_1, g_2, g_3, g_4$
\end{itemize}

\BM\ \textit{Satz E:} Die acht Polynome erzeugen die Faktorisierungsklassen
von $GF(3)[x]/(x^{26}-1)$ vollständig; ihre Kohomologieklassen
(als Charaktere der Galois-Gruppe) bilden eine $\mathbb{Q}$-Basis der
algebraischen Zykel auf $T^4/\mathbb{Z}_3$.

Beweis: Jede Hodge-Klasse auf $T^4$ liegt in einem Eigenraum des Frobenius.
Die Eigenräume werden durch die irreduziblen Faktoren des charakteristischen
Polynoms von $F_3^*$ auf $H^i$ klassifiziert. Diese Faktoren sind genau die
acht Polynome aus Dok.~342. Da jeder Eigenraum algebraisch ist (Satz~C),
sind alle Hodge-Klassen algebraisch. \hfill$\square$

% ============================================================
\section{Die Hodge-Vermutung für T⁴/Z₃}
% ============================================================

\BM\ \textit{Hauptsatz:} Alle Hodge-Klassen auf $T^4/\mathbb{Z}_3$ sind
rationale Linearkombinationen algebraischer Zykel.

Beweis aus den Sätzen A--E:
\begin{enumerate}
\item $H^*(T^4)^{\mathbb{Z}_3}$ (der ungetwistete Sektor): Alle Hodge-Klassen
      stammen aus dem $\mathbb{Z}_3$-invarianten Teil von $H^*(T^4)$.
      Da $T^4$ als abelsche Fläche die Hodge-Vermutung erfüllt (Satz~A)
      und die $\mathbb{Z}_3$-Projektion algebraische Zykel auf algebraische
      Zykel schickt, sind alle invarianten Hodge-Klassen algebraisch.
\item Twisted-Sektor: Die neun Fixpunkte sind algebraische 0-Zykel (Satz~B);
      ihre Kohomologieklassen sind trivial algebraisch.
\item Die Frobenius-Struktur (Satz~C) garantiert, dass keine Hodge-Klasse
      außerhalb dieser beiden Beiträge liegt. \hfill$\square$
\end{enumerate}

\noindent\textit{Klartext:} Für die spezifische Geometrie der FFGFT ist die
Hodge-Vermutung kein offenes Problem. Sie folgt aus der Kombination von
abelsche-Fläche-Satz, Orbifold-Kohomologie und Frobenius-Struktur. Das
Millennium-Problem in allgemeiner Form ist damit nicht gelöst.

% ============================================================
\section{Was FFGFT nicht liefert}
% ============================================================

\paragraph{Allgemeine Varietäten.}
Die Hodge-Vermutung ist für glatte projektive Varietäten beliebiger Dimension
offen. $T^4/\mathbb{Z}_3$ ist eine sehr spezielle Klasse — ein Torus-Orbifold
mit reichhaltiger algebraischer Struktur. Für eine allgemeine Varietät vom
Hodge-Typ $(2,2)$ in Dimension~4 greift kein der obigen Argumente.

\paragraph{Keine Singularitäten in der FFGFT-Lesart.}
Die FFGFT arbeitet nicht auf dem Quotientenraum $T^4/\mathbb{Z}_3$ als
Punktmenge, sondern auf dem glatten $T^4$ mit $\mathbb{Z}_3$ als
Gittersymmetrie der Moden (Dok.~314: Geometrie $\to$ Spektrum + Faser;
der Orbifold ist Gittersymmetrie, keine Zusatzannahme). Die Kohomologie ist
die $\mathbb{Z}_3$-äquivariante Kohomologie einer glatten abelschen Fläche
$E_\omega\times E_\omega$; Satz~A greift direkt, ohne Umweg über
Singularitäten. \BM

Bildet man dennoch den Quotientenraum, so hat die $\mathbb{Z}_3$-Wirkung
nach Dok.~330 die Eigenwerte $\{\omega,\omega^2\}$ auf $\mathbb{C}^2$, also
$\det=1$: Sie liegt in $SU(2)$. Die neun Fixpunkte sind dann
$A_2$-Punkte und lassen sich krepant auflösen; das Ergebnis ist eine
K3-Fläche --- glatt, projektiv, und für K3-Flächen ist die Hodge-Vermutung
bekannt (Lefschetz auf $H^{1,1}$ genügt in Dimension~2) \EM. In beiden
Lesarten liegt kein Hindernis vor.

\paragraph{Prämisse und Repräsentationsebene.}
Höherdimensionale Räume und Objekte dürfen frei konstruiert werden, solange
sie Übersetzungen des $T^4$-Inhalts sind (Typ~III, Dok.~270/330 \KM): Sie
enthalten nicht mehr als $T^4$, und keine Projektion aus ihnen kann mehr
liefern als der Modeninhalt von $T^4$ (Dok.~314: Gitter $=$ Indexmenge der
Fouriermoden; Dok.~343: die Auflösungsgrenze liegt im Gegenstand). Alle
Sätze dieses Dokuments sind solche Übersetzungen --- dieselbe kompakte
Geometrie in der Galois-Sprache von $GF(27)^*$ und in der Sprache der
Hodge-Zerlegung, bijektiv und verlustfrei. Das Dokument enthält keine
Typ-I- oder Typ-II-Operation (Dualitäts-Dekompaktifizierung, geometrische
Projektion) und trifft daher keine Aussage über die Beobachtbarkeit
einzelner Zykel; die regelt die Projektionstheorie (Dok.~285, 330, 333).

\paragraph{Nicht-kommutative Erweiterung.}
Die orbifold-Kohomologie (Chen--Ruan) stimmt im ungetwisteten Sektor mit
der klassischen Hodge-Zerlegung überein; der Twisted-Sektor hat einen
Grading-Shift, der die Hodge-Symmetrie modifiziert. Die hier gegebene
Aussage gilt für das klassische Bild; die Chen--Ruan-Erweiterung ist
als [S] einzustufen.

% ============================================================
\section{Zusammenfassung}
% ============================================================

\begin{center}
\begin{tabular}{p{4.8cm}p{4.0cm}p{2.0cm}l}
\toprule
FFGFT-Struktur & Hodge-Entsprechung & Dok. & Status \\
\midrule
$T^4$ als abelscher 2-Torus & Hodge-Vermutung bekannt & 001 & \EM \\
9 $\mathbb{Z}_3$-Fixpunkte & Algebraische 0-Zykel & 001/330 & \BM \\
Frobenius $x \mapsto x^3$ auf $GF(27)^*$ & Algebraischer Frobenius auf $H^1_{\mathrm{\acute{e}t}}$ & 339/341 & \BM \\
Sektorpaarung $k \leftrightarrow -k$ & Hodge-Symmetrie $H^{p,q} \leftrightarrow H^{q,p}$ & 343 & \BM \\
$\chi$-Klassen (8 Polynome) & Basis algebraischer Zykel & 342 & \BM \\
Hodge-Vermutung für $T^4/\mathbb{Z}_3$ & Alle Hodge-Klassen algebraisch & --- & \BM \\
Allgemeine Hodge-Vermutung & Millennium-Problem offen & --- & --- \\
\bottomrule
\end{tabular}
\end{center}

Die Hodge-Vermutung ist kein offenes Problem für $T^4/\mathbb{Z}_3$.
Die Galois-Struktur von $GF(27)^*$ — Frobenius, Sektorpaarung,
$\chi$-Klassen — ist nicht nur eine Analogie zur Hodge-Zerlegung,
sondern ihre algebraische Realisierung für diese Geometrie.
Das allgemeine Millennium-Problem bleibt unberührt.
